\section{Under Pontius Pilate}

The Christian creed contains the name of a Roman governor. It passes from the eternal generation of the Son, through the creation of all things, to a young woman of Israel and then to Pontius Pilate. The name is not there because Pilate understood what happened before him. It is there because something happened before him. The creed says that the Son ``came down from heaven'' and ``became man.'' It then says, ``For our sake he was crucified under Pontius Pilate.'' Christian confession refuses to let the event dissolve into an instructive myth or a timeless symbol.\endnote{The historical specificity belongs to the Church's confession, not merely to later apologetic embellishment. The short phrases use the current official U.S.-English Nicene Creed published by the \sourceurl{https://www.usccb.org/beliefs-and-teachings/what-we-believe}{United States Conference of Catholic Bishops}. See also \emph{Catechism of the Catholic Church} (CCC) 422--426, especially 423, which identifies Jesus of Nazareth as a Jew born of a daughter of Israel, locates his birth under Herod and Augustus, and his crucifixion under Pontius Pilate during the reign of Tiberius; \sourceurl{https://www.vatican.va/content/catechism/en/part_one/section_two/chapter_two.html}{official English}.}

That refusal can feel almost indecorous. An ultimate principle is allowed to be universal. A God who chooses Abraham, speaks through Israel, is born in Bethlehem, and gives himself under signs of water, bread, and wine can seem embarrassingly local. The more definite the claim becomes, the more it appears to some minds to have fallen beneath divine dignity. ``The transcendent'' sounds spacious; Jesus of Nazareth sounds narrow. An impersonal ultimacy makes no provincial demands and bears no marks of one people's history.

There are serious religious and philosophical traditions that use impersonal or apophatic language, and they deserve to be understood in their own terms. Nirvana, for example, is not a convenient synonym for an impersonal God, and the several Buddhist traditions cannot be reduced to a Christian polemical category. The concern here is narrower. It is the temptation, often found within a culture still borrowing Christian words, to welcome the sacred so long as it remains sacredness in general: incapable of choosing, speaking, commanding, becoming flesh, or standing before us as someone.\endnote{The comparison is deliberately bounded. Vatican II speaks of Buddhism ``in its various forms,'' recognizes its perception of the radical insufficiency of the changeable world, and insists that the Church reject nothing true and holy in other religions: Second Vatican Council, declaration \emph{Nostra aetate} (28 October 1965), no.~2; \sourceurl{https://www.vatican.va/archive/hist_councils/ii_vatican_council/documents/vat-ii_decl_19651028_nostra-aetate_en.html}{official English}. This article neither defines nirvana nor treats Buddhism as covert theism.}

The difficulty is sometimes described as resistance to a ``particular God.'' The phrase points toward something real, but taken literally it is false. God is not one member of a class called gods, as this tree is one member of a species. Nor is Christianity defending one local deity against larger rivals. What must be defended is not a particular specimen of divinity, but the freedom of the universal cause to produce particular effects, the freedom of the living God to address particular persons, and the possibility that his universal purpose should enter history by one actual road.

The creed's Roman name therefore poses a sharper question than it first appears to do. Does universality require indeterminacy? Or have we confused the living source of all things with the thinness of our most general ideas?
